%%%%%%%%%%%%%%%%%%%%%%%%%%%%%%%%%%%%%%%%%%%%%%%%%%%%%%%%%%%%%%%%%%%%%%%%%%%%%%%%
%%%%%%%%%%%%%%%%%%   Vorlage für eine Abschlussarbeit   %%%%%%%%%%%%%%%%%%%%%%%%
%%%%%%%%%%%%%%%%%%%%%%%%%%%%%%%%%%%%%%%%%%%%%%%%%%%%%%%%%%%%%%%%%%%%%%%%%%%%%%%%

% Erstellt von Maximilian Nöthe, <maximilian.noethe@tu-dortmund.de>
% ausgelegt für lualatex und Biblatex mit biber

% Kompilieren mit
% lualatex dateiname.tex
% biber dateiname.bcf
% lualatex dateiname.tex
% lualatex dateiname.tex
% oder einfach mit:
% make

\documentclass[
  tucolor,
  BCOR=12mm,     % 12mm binding corrections, adjust to fit your binding
  parskip=half,  % new paragraphs start with half line vertical space
  open=any,      % chapters start on both odd and even pages
  cleardoublepage=plain,  % no header/footer on blank pages
]{tudothesis}


% Warning, if another latex run is needed
\usepackage[aux]{rerunfilecheck}

% just list chapters and sections in the toc, not subsections or smaller
\setcounter{tocdepth}{1}

%------------------------------------------------------------------------------
%------------------------------ Sprache und Schrift: --------------------------
%------------------------------------------------------------------------------
\usepackage{fontspec}
\defaultfontfeatures{Ligatures=TeX}  % -- becomes en-dash etc.

% german language
\usepackage{polyglossia}
\setdefaultlanguage{german}

% for english abstract and english titles in the toc
\setotherlanguages{english}

% intelligent quotation marks, language and nesting sensitive
\usepackage[autostyle]{csquotes}

% microtypographical features, makes the text look nicer on the small scale
\usepackage{microtype}

%------------------------------------------------------------------------------
%------------------------ Für die Matheumgebung--------------------------------
%------------------------------------------------------------------------------

\usepackage{amsmath}
\usepackage{amssymb}
\usepackage{mathtools}
\usepackage{physics}

% Enable Unicode-Math and follow the ISO-Standards for typesetting math
\usepackage[
  math-style=ISO,
  bold-style=ISO,
  sans-style=italic,
  nabla=upright,
  partial=upright,
]{unicode-math}
\setmathfont{Latin Modern Math}

% nice, small fracs for the text with \sfrac{}{}
\usepackage{xfrac}


%------------------------------------------------------------------------------
%---------------------------- Numbers and Units -------------------------------
%------------------------------------------------------------------------------

\usepackage[
  locale=DE,
  separate-uncertainty=true,
  per-mode=symbol-or-fraction,
]{siunitx}
\sisetup{math-micro=\text{µ},text-micro=µ}

%------------------------------------------------------------------------------
%-------------------------------- tables  -------------------------------------
%------------------------------------------------------------------------------

\usepackage{booktabs}       % stellt \toprule, \midrule, \bottomrule

%------------------------------------------------------------------------------
%-------------------------------- graphics -------------------------------------
%------------------------------------------------------------------------------

\usepackage{graphicx}
\usepackage{grffile}

% allow figures to be placed in the running text by default:
\usepackage{scrhack}
\usepackage{float}
\floatplacement{figure}{htbp}
\floatplacement{table}{htbp}

% keep figures and tables in the section
\usepackage[section, below]{placeins}


%------------------------------------------------------------------------------
%---------------------- customize list environments ---------------------------
%------------------------------------------------------------------------------

\usepackage{enumitem}

%------------------------------------------------------------------------------
%------------------------------ Bibliographie ---------------------------------
%------------------------------------------------------------------------------

\usepackage[
  backend=biber,   % use modern biber backend
  autolang=hyphen, % load hyphenation rules for if language of bibentry is not
                   % german, has to be loaded with \setotherlanguages
                   % in the references.bib use langid={en} for english sources
]{biblatex}
\addbibresource{references.bib}  % die Bibliographie einbinden
\DefineBibliographyStrings{german}{andothers = {{et\,al\adddot}}}

%------------------------------------------------------------------------------
%------------------------------ Sonstiges: ------------------------------------
%------------------------------------------------------------------------------

\usepackage[pdfusetitle,unicode,linkbordercolor=tugreen]{hyperref}
\usepackage{bookmark}
\usepackage[shortcuts]{extdash}

%------------------------------------------------------------------------------
%-------------------------    Angaben zur Arbeit   ----------------------------
%------------------------------------------------------------------------------

\author{Timo Gräßer}
\title{Inverser Faraday-Effekt in Mottisolatoren}
\date{2017}
\birthplace{Lüdenscheid}
\chair{Lehrstuhl für Theoretische Physik I}
\division{Fakultät Physik}
\thesisclass{Bachelor of Science}
\submissiondate{15. August 2017}
\firstcorrector{Prof.~Dr.~Götz S. Uhrig}
\secondcorrector{Prof.~Dr.~Zweitgutachter}

% tu logo on top of the titlepage
\titlehead{\includegraphics[height=1.5cm]{logos/tu-logo.pdf}}

\begin{document}
\frontmatter
\thispagestyle{empty}
\setcounter{page}{2}
\section*{Hinweise}
Empfohlen wird die Verwendung dieser Vorlage mit der jeweils aktuellsten TeXLive Version (Linux, Windows) bzw. MacTeX Version (MacOS).
Aktuell ist dies TeXLive 2016. Download hier:
\begin{center}
  \ttfamily\url{https://www.tug.org/texlive/}
\end{center}
Bei Verwendung von TexLive Versionen 2014 und älter sollte
die Zeile
\begin{center}
\verb+\RequirePackage{fixltx2e}+ 
\end{center}
als erste Zeile der Präambel noch vor der Dokumentenklasse eingefügt werden.
Dies lädt diverse Bugfixes für LaTeX, die ab TexLive 2015 Standard sind.

Die Vorlage \texttt{thesis.tex} ist für die Kompilierung mit \texttt{lualatex} ausgelegt, mit wenigen Anpassungen kann sie aber auch mit \texttt{pdflatex} oder \texttt{xelatex} verwendet werden. 
Die Dokumentenklasse \texttt{tudothesis.cls} kann mit allen drei Programmen verwednet werden.

Achten Sie auch auf die Kodierung der Quelldateien.
Bei Verwendung von Xe\LaTeX\ oder Lua\LaTeX\ (empfohlen) müssen die
Quelldateien UTF-8 kodiert sein.
Bei Verwendung von pdf\LaTeX\ nutzen Sie die Pakete \texttt{inputenc} und \texttt{fontenc} mit der korrekten Wahl der Kodierungen.

Eine aktuelle Version dieser Vorlage steht unter 
\begin{center}
  \ttfamily\url{https://github.com/maxnoe/tudothesis}
\end{center}
zur Verfügung.

Alle verwendeten Pakete werden im \LaTeX{} Kurs von Pep et al.\ erklärt:
\begin{center}
  \ttfamily\url{http://toolbox.pep-dortmund.org/notes}
\end{center}

Für Rückmeldungen und bei Problemen mit der Klasse oder Vorlage, bitte ein \emph{Issue} auf GitHub aufmachen oder eine Email an
\href{mailto:maximilian.noethe@tu-dortmund.de}{maximilian.noethe@tu-dortmund.de} schreiben.

Wenn Sie die Dokumentenklasse mit der Option \texttt{tucolor} laden, werden verschiedene Elemente in TU-Grün gesetzt.

\maketitle

% Gutachterseite
\makecorrectorpage

% hier beginnt der Vorspann, nummeriert in römischen Zahlen
\thispagestyle{plain}

\section*{Kurzfassung}
Hier steht eine Kurzfassung der Arbeit in deutscher Sprache inklusive der Zusammenfassung der
Ergebnisse.
Zusammen mit der englischen Zusammenfassung muss sie auf diese Seite passen.

\section*{Abstract}
\begin{english}
The abstract is a short summary of the thesis in English, together with the German summary it has to fit on this page.
\end{english}

\tableofcontents

\mainmatter
% Hier beginnt der Inhalt mit Seite 1 in arabischen Ziffern
\chapter{Einleitung}
\label{ch:einleitung}

Der inverse Faraday Effekt ist ein bedeutendes Phänomen in der Magnetooptik. Er beschreibt
die Anregung einer Magnetisierung in einem Material durch die Bestrahlung mit polarisiertem Licht
und kann somit Anwendungen in technischen Bereichen finden, wie z.B. in der Daten-Speicherung und -Übertragung.
Zudem kann der Effekt zur Erzeugung von Spinwellen genutzt werden, da diese entstehen,
wenn die Magnetisierung in einem magnetisch geordneten Material aus dem Gleichgewicht gebracht wird.
In einem Experiment mit dem Titel "Magnon Accumulation by Clocked Laser Excitation as Source of
Long-Range Spin Waves in Transparent Magnetic Films"\cite{jäckl} wurden beispielsweise ferri-magnetische
Eisen-Granat-Proben mit Laserpulsen angeregt, um kohärente Spinwellen zu generieren.
Dabei diente zirkular polarisiertes Laserlicht mittels des inversen Faraday Effekts als externe periodische Kraft auf
die magnetischen Momente der Proben.
Die Verwendung von ferri-magnetischen Eisen-Granaten ist auf ihre isolierende und somit transparente Eigenschaft
zurückzuführen.  \cite{hertel, jäckl}

Im Rahmen dieser theoretischen Bachelorarbeit wird der inverse Faraday Effekt in einem Mott-Isolator modelliert und analysiert, um
Erkenntnisse für die Erzeugung einer Magnetisierung durch polarisiertes Licht, beispielsweise in Bezug auf Spinwellen-Experimente, zu erhalten.
Dazu werden zunächst theoretische Grundlagen und Konventionen zum Verständnis der physikalischen und mathematischen Inhalte
der Arbeit in Kapitel \ref{ch:theorie} dargelegt. Im darauffolgenden Kapitel \ref{ch:isolatormodelle} werden verschiedene theoretische
Isolatoren und Isolator-Modelle vorgestellt, wobei der Schwerpunkt auf dem Mott-Isolator liegt.
Kapitel \ref{ch:ife} umfasst eine kurze Herleitung zur angeregten Magnetisierung durch den inversen Faraday Effekt in Metallen und Isolatoren.
Die Modellierungen und Ergebnisse dieser Arbeit werden in Kapitel \ref{ch:ergebnisse} präsentiert. Hierbei wird mehrfach die Eignung der Modellierung
des Mott-Isolators anhand der theoretischen Erwartungen aus den Kapiteln \ref{ch:theorie} und \ref{ch:isolatormodelle} überprüft.
Darüber hinaus wird die Abhängigkeit der angeregten Magnetisierung im Mott-Isolator von der Frequenz und der Amplitude des
polarisierten Lichts analysiert.

\chapter{Theoretische Grundlagen und Konventionen}

Vektoren fett
Im Folgenden wird in natürlichen Einheiten gerechnet, das heißt unter anderem
\begin{align}
  \hbar = 1.
\end{align}
ortsraum

\section{Die Schrödingergleichung}
\label{sec:schroedingergleichung}

Die zeitliche Entwicklung eines quantenmechanischen, nicht-relativistischen Zustandes $\ket{\Psi}$ ist durch die Schrödingergleichung
\begin{align}
  \symup{i} \frac{\partial}{\partial t} \ket{\Psi(\symbf{x},t)} = H(\symbf{x},t) \ket{\Psi(\symbf{x},t)}
  \label{eqn:schroedingergleichung}
\end{align}
gegeben. Dabei ist
\begin{align}
  H(\symbf{x},t) = -\frac1{2 m} \symbf{\nabla}^2 + V(\symbf{x},t)
  \label{eqn:hamiltonian}
\end{align}
der Hamiltonoperator mit einem beliebigen Potential $V(\symbf{x},t)$.
Wenn das Potential $V$ und somit auch der Hamiltonian $H$ zeitunabhängig ist, werden in der Schrödingergleichung \eqref{eqn:schroedingergleichung} Ort $\symbf{x}$ und Zeit $t$ mit dem Ansatz
\begin{align}
  \Psi(\symbf{x},t) = \Phi(\symbf{x})f(t)
  \label{eqn:separation}
\end{align}
separiert. Die daraus resultierende stationäre Schrödingergleichung
\begin{align}
  H \ket{\Phi(\symbf{x})} = E \ket{\Phi(\symbf{x})}
  \label{eqn:statschroedinger}
\end{align}
ist nur noch vom Ort $\symbf{x}$ abhängig und die Wellenfunktionen $\Psi(\symbf{x},t)$, die die Schrödingergleichung \eqref{eqn:schroedingergleichung} lösen, sind Eigenzustände des Hamiltonoperators $H$ zu den Eigenwerten $E$.
Wenn das Potential $V$ und der Hamiltonoperator $H$ zeitabhängig sind, kann die Schrödingergleichung \eqref{eqn:schroedingergleichung} nicht separiert werden, sondern in den meisten Fällen nur numerisch, zum Beispiel durch ein Integrationsverfahren, gelöst werden.
Nähere Informationen zur Herleitung und zur Lösung der Schrödingergleichung können \cite{schwabl} entnommen werden.

\section{Quantenmechanische Operatoren und die zweite Quantisierung}

In der Quantenmechanik sind den Observablen der klassischen Physik Operatoren zugeordnet.
Dieses Postulat wird als Korrespondenzprinzip bezeichnet.
\cite{schwabl}

Die zweite Quantisierung ist eine Darstellung von Operatoren, die sich vor allem für die Betrachtung von fermionischen Vielteilchensystemen eignet.
Grundlage für diese Darstellung sind die auf dem Fockraum definierten fermionischen Erzeugungs- und Vernichtungsoperatoren $c^\dag$ und $c$.
Diese ändern die Teilchenzahl eines fermionischen Vielteilchenzustands gemäß
\begin{align}
  c_i^{\phantom{\dag}} \ket{N_1, N_2, ..., N_i, ...} & = \pm n_i \ket{N_1, N_2, ..., N_i-1, ...} \label{eqn:vernichter}\\
  c_i^\dag \ket{N_1, N_2, ..., N_i, ...} & = \pm (1-N_i) \ket{N_1, N_2, ..., N_i+1, ...} \label{eqn:erzeuger}
\end{align}
um $1$. Dabei ist $N_i$ die Besetzungszahl des i-ten Zustands, die aufgrund des Pauliprinzips nur die Werte $0$ und $1$ annehmen kann.
Das Vorzeichen in \eqref{eqn:vernichter} und \eqref{eqn:erzeuger} entsteht durch die Antisymmetrie des Zustands. \cite{cy}

Im Folgenden werden Operatoren und Zustände ausschließlich in Ortsdarstellung angegeben und verwendet.

\section{Der Besetzungszahloperator}

Alle fermionischen Vielteilchenzustände sind Eigenfunktionen der Besetzungszahloperatoren
\begin{align}
  n_i = c_i^\dag c_i^{\phantom{\dag}},
  \label{eqn:besetzer}
\end{align}
wobei die zugehörigen Eigenwerte jeweils die Besetzungszahlen $n_i$ der i-ten Zustände sind.
Der Wertebereich der Eigenwerte ist daher ebenfalls auf $0$ und $1$ beschränkt.
\cite{cy}

\section{Der Spinoperator}

Die Indizes der Erzeugungs- und Vernichtungsoperatoren werden im Folgenden um einen zusätzlichen Parameter $\sigma$, der die Spinausrichtung eines Elektrons angibt, erweitert.
Die Ein-Teilchen-Zustände des fermionischen Vielteilchensystems spalten sich dadurch jeweils in zwei im Ort entartete Zustände auf, die sich nur in der Spin-Magnetquantenzahl $\sigma$ unterscheiden.
Der Parameter $\sigma$ kann die Werte $\sigma = \frac12$, für Spin Up ($\uparrow$), und $\sigma = -\frac12$, für Spin Down ($\downarrow$), annehmen.

Um den Gesamtspin $S$ eines fermionischen Vielteilchensystems zu ermitteln, wird zunächst der vektorartige Spinoperator für das Einteilchensystem
\begin{align}
  S & = \begin{pmatrix*}[r]
    S_x \\
    S_y \\
    S_z
  \end{pmatrix*}
  \label{eqn:spinoperator}
\end{align}
eingeführt. Die Quantisierungsachse wird im Folgenden in z-Richtung gelegt, sodass die Einteilchenzustände Eigenfunktionen des $S_z$-Operators zu den Eigenwerten $\sigma$ sind.
Der $S_x$- und der $S_y$-Operator werden dann mit
\begin{align}
  S_x & = \frac12 \left(S_+ + S_- \right) & S_y &= \frac1{2i} \left(S_+ - S_- \right)
  \label{eqn:leiteroperatorenausdruck}
\end{align}
durch die Spin-Leiteroperatoren $S_+$ und $S_-$ ausgedrückt. Die Wirkung der Spin-Leiteroperatoren auf einen fermionischen Einteilchenzustand mit Spin Up oder Spin Down ist durch
\begin{align}
  S_+ \ket{\uparrow} & = 0 & S_+ \ket{\downarrow} & = \ket{\uparrow} \label{eqn:spluswirkung}\\
  S_- \ket{\uparrow} & = \ket{\downarrow} & S_- \ket{\downarrow} & = 0 \label{eqn:sminuswirkung}
\end{align}
gegeben. Das Anwenden des $S_+$-Operators auf einen Einteilchenzustand entspicht dem Vernichten eines Elektrons mit Spin Down gefolgt vom Erzeugen
eines Elektrons mit Spin Up. Dies gilt auch für den $S_-$-Operator mit vertauschten Spins. Daraus folgt Darstellung der Spin-Leiteroperatoren in zweiter Quantisierung
\begin{align}
  S_+ & = c_\uparrow^\dag c_\downarrow^{\phantom{\dag}} & S_- & = c_\downarrow^\dag c_\uparrow^{\phantom{\dag}}. \label{eqn:leitersecondquant}
\end{align}
Die Eigenwertgleichungen für den $S_z$-Operator
\begin{align}
  S_z \ket{\uparrow} &= \frac12 \ket{\uparrow} & S_z \ket{\downarrow} &= -\frac12 \ket{\downarrow}
\end{align}
führen auf die Darstellung des Operators in zweiter Quantisierung
\begin{align}
  S_z = \frac12 (n_\uparrow -  n_\downarrow) = \frac12 (c_\uparrow^\dag c_\uparrow^{\phantom{\dag}} - c_\downarrow^\dag c_\downarrow^{\phantom{\dag}}) \label{eqn:szsecondquant}.
\end{align}
\cite{schwabl}

Aus den Gleichungen \eqref{eqn:spinoperator} und \eqref{eqn:leiteroperatorenausdruck} folgt der $S^2$-Operator für ein fermionisches Einteilchensystem
\begin{align}
  S^2 & = S_z^2 + \frac12 \left( S_+ S_- + S_- S_+ \right)
  \label{eqn:squadoperatorein}
\end{align}
Durch summieren der Einteilchenoperatoren $S_{+,i}$, $S_{-,i}$ und $S_{z,i}$ über alle Zustandsniveaus $i$ folgt aus den Gleichungen \eqref{eqn:squadoperatorein} die Darstellung des $S^2$-Operators in zweiter Quantisierung für ein fermionisches Vielteilchensystem
\begin{align}
  S^2 & = \frac14 \left(\sum_{i} n_{\uparrow,i} - n_{\downarrow,i}\right)^2 +
  \frac12 \left( \sum_{i,j} c_{\uparrow,i}^\dag c_{\downarrow,i}^{\phantom{\dag}} \cdot c_{\downarrow,j}^\dag c_{\uparrow,j}^{\phantom{\dag}} + c_{\downarrow,i}^\dag c_{\uparrow,i}^{\phantom{\dag}} \cdot c_{\uparrow,j}^\dag c_{\downarrow,j}^{\phantom{\dag}} \right).
  \label{eqn:squadoperatorviel}
\end{align}
Die Vielteilchenzustände sind Eigenfunktionen des $S^2$-Operators zu den Eigenwerten $S(S+1)$.
\cite{???}

\chapter{Festkörper-Modelle}

evtl kurz Bändermodell

\section{Das Tight-binding-Modell und der Stromoperator}

Das allgemeine Bändermodell beruht auf der Annahme, dass sich die Elektronen in einem Festkörper unter dem Einfluss eines periodischen Gitterpotentials quasifrei bewegen.
Durch das Tight-Binding-Modell wird ein Festkörpersystem beschrieben, in welchem die Elektronen stark gebunden und lokalisiert sind und abgesehen von der Pauliabstoßung nicht miteinander in Wechselwirkung treten.
Für die Elektronen sind dann nur noch diskrete Sprünge zwischen den Gitterplätzen des Festkörpers möglich. Der Hamiltonoperator des Tight-Binding-Modells in zweiter Quantisierung lautet
\begin{align}
  H_\text{tb} = -J \sum_{i,\sigma} (c_{i+1,\sigma}^{\dag}c_{i,\sigma}^{\phantom{\dag}} + c_{i,\sigma}^{\dag}c_{i+1,\sigma}^{\phantom{\dag}})
  \label{eqn:hamiltontb}
\end{align}
für ein Festkörpersystem mit einem Elektronenband. Summiert wird dabei über alle Gitterpätze $i$ und Spinausrichtungen $\sigma$.
Der Term $c_{i+1,\sigma}^{\dag}c_{i,\sigma}^{\phantom{\dag}}$ lässt ein Elektron vom Gitterplatz $i$ zu $i+1$ springen und der Term $c_{i,\sigma}^{\dag}c_{i+1,\sigma}^{\phantom{\dag}}$ umgekehrt.
Der Vorfaktor $J$ ist eine energetische Konstante, welche die Stärke der Tight-Binding-Wechselwirkung im Hamiltonian beschreibt.
Stromoperator!!!!!!!!!!!!

\section{Das Hubbard-Modell}

Die bisher erwähnten Modelle vernachlässigen die Coulomb-Wechselwirkung der Elektronen untereinander.
Das Hubbard-Modell ist eine spezielle Erweiterung des Tight-Binding-modells, welche diese Wechselwirkung berücksichtigt.
Dies ist im Hamiltonoperator durch einen zusätzlichen Term, der eine starke Coulombabstoßung zwischen den Elektronen beschreibt, realisiert.
Im Folgenden wird die Abstoßung in erster Ordnung, also für Elektronen, die sich auf demselben Gitterpatz befinden, betrachtet.
Diese Näherung lässt sich dadurch begründen, dass das Coulombpotential mit $\frac1{r}$ abfällt. Der Hamiltonoperator des Hubbardmodells ist dann in zweiter Quantisierung durch
\begin{align}
  H_\text{hubb} = H_\text{tb} + U\sum_{i} \prod_{\sigma} n_{i,\sigma}.
  \label{eqn:hamiltonhubb}
\end{align}
gegeben, wobei erneut über die Gitterplätze $i$ und die Spinausrichtungen $\sigma$ summiert wird.

\section{Der Übergang zum Heisenberg-Austausch-Modell im Grenzfall des halbgefüllten Hubbard-Modells}

Die Besetzungswahrscheinlichkeit der Zustände im halbgefüllten Hubbardmodell, in denen sich jeweils ein Elektron auf jedem Gitterplatz befindet, steigt mit dem Wechselwirkungs-Potential $U$ an, da
die Doppelbesetzung eines Atoms mit dem Ansteigen des Potentials $U$ energetisch ungünstiger wird. Im Grenzfall $U \to \infty$ besteht der Hilbertraum des Festkörpersystems nur noch aus diesen Zuständen
und die Elektronen können sich nicht mehr zwischen den Gitterplätzen bewegen. Die Dynamik des Systems ist dann nur noch durch virtuelle Prozesse gegeben, in welchen die Spinausrichtung der Elektronen
benachbarter Gitterplätze getauscht oder beibehalten wird. Für einen virtuellen Prozesse zwischen zwei benachbarten Elektronen mit antiparalleler Spinausrichtung, z.B.
\begin{align}
  (\uparrow; \downarrow) &\to (\uparrow \downarrow ; 0) \to (\uparrow; \downarrow),
\end{align}
ergibt sich in zweiter Ordnung Störungstheorie die Energiedifferenz
\begin{align}
  \Delta E_{\uparrow \downarrow} = - \frac{2J^2}{U}.
  \label{eqn:Ediffantipar}
\end{align}
Ein Sprung zwischen zwei benachbarten Elektronen mit paralleler Spinausrichtung ist wegen des Pauliprinzips nicht möglich.
Der Energieunterschied beträgt daher
\begin{align}
  \Delta E_{\uparrow \uparrow} = 0.
  \label{eqn:Ediffpar}
\end{align}
Da die virtuellen Prozesse nur noch von der Spinausrichtung abhängen, wird das System mit dem Heisenberg-Austausch-Modell verglichen.
Der Hamiltonoperator dieses Modells ist durch
\begin{align}
  H_\text{spin} = t \sum_{i,j} (S_i S_j + C)
  \label{eqn:hamiltonspincj}
\end{align}
gegeben, wobei in beiden Summen über alle Gitterplätze $i$ und $j$ summiert wird. Durch das Vergleichen der Energiedifferenzen der unterschiedlichen Spinkonfigurationen in diesem Modell
\begin{align}
  \Delta E_{\uparrow \downarrow} = \bra{\uparrow \downarrow} H_\text{spin} \ket{\uparrow \downarrow} = t\left(C - \frac14\right) \\
  \Delta E_{\uparrow \uparrow} = \bra{\uparrow \uparrow} H_\text{spin} \ket{\uparrow \uparrow} = t\left(C + \frac14\right)
\end{align}
mit den Energiedifferenzen \eqref{eqn:Ediffantipar} und \eqref{eqn:Ediffpar} ergeben sich die Konstanten
\begin{align}
  t = \frac{4 J^2}{U}
\end{align}
und
\begin{align}
  C = -\frac14
\end{align}
für den Hamiltonian \eqref{eqn:hamiltonspincj}.
Im Grenzfall $U \to \infty$ geht der Hamilton-Operator des halbgefüllten Hubbard-Modells \eqref{eqn:hamiltonhubb} also in den Hamilton-Operator des Heisenberg-Austausch-Modells
\begin{align}
  H_\text{spin} = \frac{4 J^2}{U} \sum_{i,j} (S_i S_j - \frac14)
  \label{eqn:hamiltonspin}
\end{align}
über.

Eine mögliche Struktur der Arbeit sieht wie folgt aus:

\begin{enumerate}
    \item \textbf{Einleitung}\\
        In der \emph{kurzen} Einleitung wird die Motivation für die Arbeit
        dargestellt und ein Einblick in die kommenden Kapitel gegeben.
    \item \textbf{Theoretische Grundlagen}\\
        Alles was an theoretischen Grundlagen benötigt wird, sollte auch eher kurz gehalten werden.
        Statt Grundlagenwissen zu präsentieren, eher auf die entsprechenden Lehrbücher verweisen.
        Etwa: Tiefer gehende Informationen zur klassischen Mechanik entnehmen Sie bitte \cite{kuypers}.
    \item \textbf{Ergebnisse} \\
        Der eigentliche Teil der Arbeit, das was getan wurde.
    \item \textbf{Zusammenfassung und Ausblick} \\
        Zusammenfassung der Ergebnisse, Optimierungsmöglichkeiten, mögliche weitergehende Untersuchungen.
\end{enumerate}

Die Gliederung sollte auf der einen Seite nicht zu fein sein, auf der anderen Seite
sollten sich klar unterscheidende Abschnitte auch kenntlich gemacht werden.

In der hier verwendeten \KOMAScript-Klasse \texttt{scrbook} ist die oberste Gliederungsebene,
die in der Bachelorarbeit verwendet werden sollte, das \texttt{\textbackslash chapter}.

Ein Kapitel sollte erst dann in tiefere Gliederungsebenen unterteilt werden, wenn es auch wirklich etwas zu unterteilen gibt. Es sollte keine Kapitel mit nur einem Unterkapitel (\texttt{\textbackslash section}) geben.

In dieser Vorlage ist die Tiefe des Inhaltsverzeichnisses auf \texttt{chapter} und \texttt{section} beschränkt. Möchten Sie diese Beschränkung aufheben, entfernen Sie den Befehl
\begin{verbatim}
            \setcounter{tocdepth}{1}
\end{verbatim}
aus der Präambel oder ändern Sie den Zahlenwert entsprechend. Das Inhaltsverzeichnis sollte für eine Bachelorarbeit auf eine Seite passen.

\chapter{Isolatoren}\label{make}

Diese Vorlage ist auf die Kompilierung mit \texttt{lualatex} ausgelegt.
Als Dokumentenklasse  wird die \KOMAScript\-Klasse \texttt{scrbook} verwendet.
Falls Sie Änderungen am Layout vornehmen möchten, lesen Sie die \KOMAScript-Dokumentation: \cite{koma}.

Eine umfangreiche Einführung in die moderne Verwendung von \LaTeX{} gibt es hier: \cite{toolbox}, lesenswert ist außerdem das \LaTeX-Tabu: \cite{l2tabu}

Um dieses Dokument vollständig zu erstellen sind maximal vier Programmläufe nötig:
\begin{enumerate}[nosep]
    \item \texttt{lualatex BachelorArbeit.tex}
    \item \texttt{biber BachelorArbeit.bcf}
    \item \texttt{lualatex BachelorArbeit.tex}
    \item \texttt{lualatex BachelorArbeit.tex}
\end{enumerate}

Beim ersten Lauf des \LaTeX-Compilers werden die Kapitel, Links und zitierten Bibliographieeinträge in Hilfsdateien geschrieben.

Dann ist ein Lauf des Programms \texttt{biber} nötig, welches die benötigten Einträge aus der Hilfsdatei einliest, die Einträge aus der \texttt{.bib} Datei einliest, sortiert und formatiert und in eine weitere Hilfsdatei schreibt.

Beim nächsten \LaTeX-Lauf werden dann diese Hilfsdateien eingelesen und Literatur- und Inhaltsverzeichnis erstellt.

Manchmal ist ein vierter Lauf nötig, falls sich durch das einfügen des Literaturverzeichnisses Seitenzahlen verändert haben.

Das Tool \texttt{latexmk} übernimmt dies mit nur einem Programmaufruf und
führ nur so viele Aufrufe durch, wie nötig sind.

\texttt{latexmk --lualatex BachelorArbeit.tex}

Eine gute Option ist es, den \LaTeX{} Output in einem anderen
Ordner zu erzeugen, dies ist mit der \texttt{--output-directory} Option möglich:

\texttt{latexmk --output-directory=build --lualatex BachelorArbeit.tex}


\section{Erstellen des Ausgabedokuments mit Make}

Für diese Vorlage wird ein Makefile zur Verfügung gestellt, welches automatisch alle Schritte ausführt, die für das fertige Dokument nötig sind.
Die Ausgabe erfolgt dabei in den Unterordner \texttt{build/}.
Make prüft, ob die Quelldateien verändert wurden, falls nicht, werden auch keine Befehle ausgeführt.

Falls Sie das Makefile benutzen möchten, sollten Sie alle Abhängigkeiten eintragen (Eigene Dateien für Kapitel, Plots, etc.).


Download und weitere Informationen zu Make gibt es unter \cite{make}. Die Befehle sind für die Bash ausgelegt.
Wenn Sie sie unter Windows nutzen wollen, benötigen Sie einen Bash-Emulator, wie Git Bash, Download unter \cite{gitbash} möglich.
Wenn Sie Make installiert haben, rufen Sie einfach in der Konsole im Verzeichnis der Arbeit den Befehl \texttt{make}.

\section{Erstellen des Ausgabedokuments mit Texmaker}
\subsection{Einrichten der nötigen Befehle}
Ein beliebter Editor für alle Betriebssysteme ist Texmaker, Download unter \cite{texmaker}.
Damit Texmaker das Dokument korrekt kompiliert, fügen sie einen benutzerdefinierten Befehl hinzu:
\begin{enumerate}[nosep]
    \item Klicken sie oben in der Menüleiste auf \emph{Benutzer/in}
    \item Klick auf \emph{Eigene Befehle}
    \item Klich auf \emph{Eigene Befehle editieren}, dort können Sie bis zu 5 eigene Befehle definieren
    \item Geben Sie dem Befehl unter \emph{Menüeintrag} einen Namen und tragen sie folgende Befehle in das Befehlsfeld ein: \\
      \verb+latexmk --lualatex --interaction=batchmode --halt-on-error %.tex |+ \\
    \item Bestätigen Sie mit \emph{OK}
\end{enumerate}

In Abbildung \ref{fig:texmaker} ist ein Screenshot des Befehlsmenü gezeigt. Ihren Befehl können Sie nun im Drop-Down-Menü zum
Kompilieren des Dokuments auswählen und mit einem Klick auf den Pfeil starten.

\subsection{Aufräumen}

Nach einem \LaTeX-Fehler ist es oft notwendig, die erstellten Hilfsdateien zu löschen.
Klicken Sie hierzu auf \emph{Werkzeuge}→\emph{Aufräumen}.


\chapter{\LaTeX-Grundlagen}

Bitte beachten Sie beim Schreiben der Arbeit folgende Konventionen bzw. Grundlagen:

\begin{itemize}
    \item \textbf{Abschnitte und Zeilenumbrüche} \\
        Es sollten im Fließtext keine Zeilenumbrüche mit \textbackslash\textbackslash \ erzwungen werden.
        Schreiben Sie höchsten einen Satz in eine Code-Zeile.
        Absätze werden im Code mit einer Leerzeile markiert und dann entsprechend der Einstellung von \texttt{parskip} in der Dokumentenklasse gesetzt.
    \item \textbf{Kursiv/Aufrecht} \\
        \begin{itemize}
            \item Variablen und physikalische Größen werden kursiv gesetzt.
            \item Einheiten werden immer aufrecht und mit einem halben Leerzeichen Abstand zur Zahl gesetzt. Nutzen Sie \texttt{siunitx}!
            \item Mathematische Konstanten und Funktionen werden ebenfalls aufrecht gesetzt. Zum Beispiel die Eulersche Zahl e, das imaginäre i und das infinitesimale d.
                Im Mathematikmodus können Sie dies mit dem Befehl \verb_\mathrm{}_ erreichen. Für die Funktionen stellt \LaTeX \ Befehle bereit, z.B. \verb+\arccos+.
            \item Integrand und ein $\mathrm{d}x$ sollten ebenfalls durch ein kleines Leerzeichen (\verb+\,+) getrennt werden.
        \end{itemize}



\end{itemize}

\section{Zahlen und Einheiten}

Jede Zahl, jede Einheit und jede Zahl mit Einheit sollte mit Hilfe der in dem Paket \texttt{siunitx} zur Verfügung gestellten Befehle gesetzt werden.
Grundsätzlich gilt: Einheiten werden aufrecht gesetzt und haben ein kleines Leerzeichen (\verb+\,+) Abstand zu ihrer Zahl.
Werden Fließkommazahlen ohne \texttt{siunitx} gesetzt, entsteht ein hässlicher Leerraum zwischen Komma und erster Nachkommastelle, da \LaTeX \ das Komma nicht als Dezimaltrennzeichen, sondern als Satzzeichen interpretiert.

Das Paket wurde mit deutschen Spracheinstellungen (also mit Komma als Dezimaltrennzeichen und $\cdot$ zwischen Zahl und Zehnerpotenz) geladen, sowie mit den Einstellungen, dass die Standardabweichung stets durch $\pm$ abgetrennt wird und Einheiten falls nötig als Brüche ausgegeben werden.

\begin{table}
    \centering
    \caption{Beispiele für siunitx}
    \label{tab:si}
    \begin{tabular}{l r}
        \toprule
        Befehl     &   Ergebnis \\
        \midrule
        \verb+\num{1.2345}+ & \num{1.2345} \\
        \verb+\num{1.2e3}+ & \num{1.2e3} \\
        \verb_\num{1.2 +- 0.2}_ & \num{1.2+-0.2} \\
        \verb+\num{10000}+ & \num{10000} \\
        \verb+\si{\meter\per\second}+ & \si{\meter\per\second} \\
        \verb+\SI{1.2(1)}{\micro\ampere}+ & \SI{1.2(1)}{\micro\ampere} \\
        \verb+\SI{1.2\pm0.1e3}{\kilo\gram\per\cubic\meter}+ & \SI{1.2\pm0.1e3}{\kilo\gram\per\cubic\meter} \\
        \bottomrule
    \end{tabular}
\end{table}

Das Paket stellt unter anderem die drei wichtigen Befehle
\begin{itemize}
    \item \texttt{\textbackslash num\{Zahl\}},
    \item \texttt{\textbackslash si\{Einheit\}} und
    \item \texttt{\textbackslash SI\{Zahl\}\{Einheit\}}
\end{itemize}
zur Verfügung.
Diese Befehle sollten stets genutzt werden, wenn Zahlen angegeben werden.
Sie funktionieren sowohl im Text- als auch im Mathematikmodus.
In Tabelle \ref{tab:si} sind einige Beispiele aufgetragen. Bitte lesen Sie die Dokumentation \cite{siunitx}.

\section{Das Literaturverzeichnis}

Das Literaturverzeichnis wird mit Hilfe von BibLaTeX und biber erstellt.
Tragen Sie alle ihre Quellen in die Datei \texttt{references.bib} ein, Sie enthält bereits
einige Beispiele. Für weitere Informationen lesen Sie bitte die Dokumentation \cite{biblatex}.

Im Text können Sie mit \verb_\cite{kürzel}_ zitieren. Seitenzahlen geben Sie in eckigen Klammern an:
\verb_\cite[10]{kürzel}_.

Das Literaturverzeichnis ist so eingestellt, dass es Ihre Quellen in alphabetischer Reihenfolge nach Autoren nummeriert.
Möchten Sie das Literaturverzeichnis nach der Reihenfolge des Auftauchens im Text sortieren, fügen sie die Paktetoption \texttt{sorting=none} beim Laden
des BibLaTeX-Pakets hinzu.

Den Zitier- und Bibliographie-Stil geben sie mit der Option \texttt{style=Stil} an. Die beiden gebräuchlisten Stile sind \texttt{numeric} und \texttt{alphabetic}.
Bei \texttt{numeric} werden die Quellen durchnummeriert, bei \texttt{alphabetic} wird ein Buchstabenkürzel aus Autor(en)-Name(n) und Jahr verwendet.
Für weitere Stile konsultieren Sie bitte die Dokumentation: \cite{biblatex}.

Ein Beispiel für das Zitieren eines Buches lautet so \cite{handbook_adhesives},
wissenschaftliche Artikel hingegen werden so \cite{einstein} zitiert.

Damit das Literaturverzeichnis erstellt wird, ist ein Aufruf von \texttt{biber} nach einem ersten kompilieren mit \texttt{lualatex} nötig.
Danach muss das Dokument erneut mit \texttt{lualatex} kompiliert werden.

Zum korrekten Kompilieren des Dokuments siehe Kapitel \ref{make}.

\chapter{Inverser Faraday-Effekt}

\section{Abbildungen}

Achten Sie bei ihren Plots auf ausreichend große Achsenbschriftungen, ausreichende Schriftdicken und gut unterscheidbare Farben.
Im Idealfall haben Sie im Plot und der Arbeit die gleiche Schriftgröße und Schriftart.
Dies lässt sich durch Erstellen des Plots in der korrekten Größe und Einbinden mit dem optionalen Argument \texttt{scale=1} erreichen. Ein Beispiel sehen Sie in Abbildung \ref{fig:bsp}.

Nutzen Sie wenn möglich Vektorgrafiken (pdf) und nur in Ausnahmen Rastergrafiken wie .png oder .jpg.
Setzen Sie Punkte hinter Abbildungsunterschriften.

\section{Tabellen}

Tabellen sollten so einfach wie möglich aufgebaut sein, verzichten Sie auf zu viele Linien. In fast allen Fällen reichen drei horizontale Linien aus, jeweils über und unter der Tabelle und zwischen den Spaltenüberschriften und der eigentlichen Tabelle.

Das Paket \texttt{booktabs} stellt hierfür \verb_\toprule_, \verb_\midrule_ und
\verb_\bottomrule_ zur Verfügung.
Das Paket \texttt{siunitx} stellt eine extrem mächtige neue Spalteneinstellung bereit: \texttt{S}, mit ihr können Zahlen und Einheiten sehr sauber und gut ausgerichtet gesetzt werden.

Diese Vorlage geht von Tabellenüberschriften aus, möchten Sie dagegen Tabellenunterschriften entfernen Sie das entsprechende optionale Argument für die Dokumentenklasse in der Präambel.

Ein Beispiel ist Tabelle~\ref{tab:bsp}.
\begin{table}
    \centering
    \caption{Beispieltabelle mit willkürlichen Werten, für die Zahlenwerte wurde die S-Option aus \texttt{siunitx} verwendet.}
    \label{tab:bsp}
    \begin{tabular}{S[table-format=4.2] S[table-format=3.2]}
        \toprule
        {$p \mathrel{/} \si{\pascal}$}  & {$T \mathrel{/} \si{\kelvin}$} \\
        \midrule
        1024,23 & 273,15 \\
        1025,31 & 274,5 \\
        1026,27 & 276,2 \\
        \bottomrule
    \end{tabular}
\end{table}

\chapter{Ergebnisse}

\section{Untersuchung eines Festkörpersystems im Tight-Binding-Modell}

Im Folgenden wird die Dynamik eines mikroskopischen Festkörpersystems, bestehend aus vier Gitterplätzen bzw. Niveaus und einem Elektron, im Tight-Binding-Modell analysiert.
Das System ist in Abbildung \ref{fig:system} mit beispielhafter Elektronenbesetzung skizziert. Das Elektron kann waagerecht oder senkrecht zwischen den Gitterplätzen springen.
Der Hilbertraum des Systems setzt sich aus vier Zuständen, bei denen sich das Elektron jeweils auf einem der vier Gitterplätze befindet, zusammen und ist daher vierdimensional.
Um die Dynamik des Systems zu analysieren werden die Zustände aufgestellt, daraus die vierdimensionale Hamilton-Matrix bestimmt und mit Hilfe dieser die Schrödingergleichung \ref{eqn:schroedinger} gelöst.

Die Zustände werden in der Konvention
\begin{align}
  \ket{i} = \ket{f_i(x)}
  \label{eqn:tbzustandsvorschrift1e}
\end{align}
dargestellt. Dabei gibt die Funktion $f_i(x)$ die Nummer des Gitterplatzes an, auf dem sich das Elektron $x$ befindet. Aus der Gleichung \eqref{eqn:hamiltontb} folgt der Hamiltonian des Systems
\begin{align}
  H_\text{tb,sys} = -J \sum_{i=1}^{4} (c_{i+1}^{\dag}c_{i} + c_{i}^{\dag}c_{i+1}).
\end{align}
Da das Elektron auch zwischen den Gitterplätzen $1$ und $4$ springen kann, sind die Randbedingungen periodisch, sodass die Indizes $i+1$ für $i=4$ den Wert $1$ annehmen.
In den folgenden Rechnungen wird der Vorfaktor $J$ auf $1$ gesetzt. Der Hamiltonian wird mit der Formel
\begin{align}
  H_{mn} = \bra{m} H \ket{n}
  \label{eqn:matrixdarstellung}
\end{align}
als $4 \times 4$ -Matrix dargestellt. Die berechnete Matrix ist in Tabelle \ref{tab:tbhamiltonmatrixe1} zu sehen.
Als nächstes wird die Schrödingergleichung \eqref{eqn:schroedinger}, die durch den Hamilton-Operator $H$ festgelegt ist, gelöst.
Da dieser zeitunabhängig ist, wird die stationäre Schrödingergleichung \eqref{eqn:statschroedinger} aufgestellt und durch Diagonalisieren des Hamilton-Operators $H$ mit Papier und Bleistift gelöst. (analytisch)
Die berechneten Energie-Eigenwerte $E_k/J$ und die Koeffizienten $\alpha_{i,k}$ der generierten Eigenzustände $v_k$ sind in Tabelle \ref{tab:tbeigene1} aufgelistet.
Die Eigenzustände ergeben sich aus der Linearkombination
\begin{align}
  v_k = \sum_i \alpha_{i,k} \ket{i}.
  \label{eqn:linkomb}
\end{align}
Da das Betragsquadrat des Koeffizienten $\alpha_{i,k}$ der Besetzungswahrscheinlichkeit $P_i$ für den i-ten Zustand im Eigenzustand $v_k$ entspricht,
ist an die Linearkombination \eqref{eqn:linkomb} die Normierungsbedingung
\begin{align}
  \sum_i \lvert \alpha_{i,k} \rvert^2 = 1.
  \label{eqn:tbnb}
\end{align}
geknüpft.

\begin{table}[h]
  \centering
  \caption{Analytisch berechnete Eigenwerte $E_k/J$ und zugehörige Koeffizienten $\alpha_{i,k}/0.5$ der Eigenzustände.}
  \begin{tabular}{S[table-format=1.0] S[table-format=2.0] S[table-format=1.0] S[table-format=1.0] S[table-format=1.0] S[table-format=2.0]}
    \toprule
    {$k$} & {$E_k/J$} & {$\alpha_{1,k}/0.5$} & {$\alpha_{2,k}/0.5$} & {$\alpha_{3,k}/0.5$} & {$\alpha_{4,k}/0.5$}\\
    \midrule
    0 & -2 & 1 & $\sqrt{2}$         & 0           & 1  \\
    1 & 0  & 1 & 0                  & $\sqrt{2}$  & -1 \\
    2 & 0  & 1 & \text{$-\sqrt{2}$} & 0           & 1  \\
    3 & 2  & 1 & 0                  & \text{$-\sqrt{2}$} & -1 \\
    \bottomrule
  \end{tabular}
  \label{tab:tbeigene1}
\end{table}

Das System ist invariant unter einer Transformation, die das Gitter um $\SI{90}{\degree}$ dreht und die Gitterplätze in derselben Konvention wie im unrotierten System neu durchnummeriert.
Da eine Rotation die Dynamik des Systems nicht ändert, kann folglich die Konvention der Durchnummerierung die Dynamik ebenfalls nicht ändern.
Die Zustände $\ket{i}$ unterscheiden sich neben der Nummer $i$ um keine weitere Quantenzahl. Daraus folgt, dass die Besetzungswahrscheinlichkeiten in den
Eigenzuständen, die sich aus nicht-entarteten Energie-Eigenwerten ergeben, für alle Zustände $\ket{i}$ gleich sind.
Aus der Normierungsbedingung \eqref{eqn:tbnb} folgt daher für diese Eigenzustände
\begin{align}
  P_{i,k} = \lvert \alpha_{i,k} \rvert^2 = 0.25.
  \label{eqn:BesetzungRot}
\end{align}
Das ist für die berechneten Eigenzustände $v_0$ und $v_3$ der Fall. Für die Eigenzustände zu entarteten Energie-Eigenwerten, hier $v_1$ und $v_2$ zur Energie $E_1 = E_2 = 0$, können mehrere verschiedene Basen
gewählt werden, sodass Gleichung \eqref{eqn:BesetzungRot} nicht zwangsläufig gilt. Der physikalische Grund dafür ist, dass das Elektron bei hinreichender Energie einen Zustand,
der sich aus einer Linearkombination der beiden Eigenzustände $v_1$ und $v_2$ ergibt, besetzt.

In dem bisher betrachteten System kann unter Beachtung des Pauliprinzips ein weiteres Elektron ergänzt werden.
Die waagerechten oder senkrechten Sprünge der Elektronen sind  zwichen den Gitterplätzen springen,




\begin{table}[h]
  \centering
  \caption{Datenpunkte der drei niedrigsten Eigenenergien des Festkörpersystems in Abhängigkeit vom Potential und in Einheiten von $J$.}
  \begin{tabular}{S[table-format=1.0] S[table-format=2.0]}
    \toprule
    {$k$} & {$E/J$}\\
    \midrule
    0  & -2 \\
    1  & -2 \\
    2  & 0  \\
    3  & 0  \\
    4  & 2  \\
    5  & 2  \\
    \bottomrule
  \end{tabular}
  \label{tab:tbeigenwertee1}
\end{table}

\begin{table}[h]
  \centering
  \caption{Datenpunkte der drei niedrigsten Eigenenergien des Festkörpersystems in Abhängigkeit vom Potential und in Einheiten von $J$.}
  \begin{tabular}{S[table-format=1.0] S[table-format=2.0]}
    \toprule
    {$k$} & {$E/J$}\\
    \midrule
    1  & -2 \\
    2  & -2 \\
    3  & 0  \\
    4  & 0  \\
    5  & 2  \\
    6  & 2  \\
    \bottomrule
  \end{tabular}
  \label{tab:tbeigenwerte}
\end{table}




In Tabelle \ref{tab:tbzustandsvorschrift} sind die Darstellungen \eqref{eqn:tbzustandsvorschrift} der vier Zustände aufgelistet.

\begin{table}
  \centering
  \caption{.}
  \begin{tabular}{S[table-format=1.0] S}
    \toprule
    {$i$} & {$\ket{i}$} \\
    \midrule
    1 & $\ket{1,2}$ \\
    2 & $\ket{1,3}$ \\
    3 & $\ket{1,4}$ \\
    4 & $\ket{2,3}$ \\
    5 & $\ket{2,4}$ \\
    6 & $\ket{3,4}$ \\
    \midrule
  \end{tabular}
  \label{tab:hubbzustandsvorschrift}
\end{table}


\section{Untersuchung eines Festkörpersystems im Hubbard-Modell}

Diese Konvention führt automatisch dazu, dass der zweite Eintrag in der Darstellung der Zustände \eqref{eqn:tbzustandsvorschrift} größer als der erste ist.
Falls der zweite Eintrag eines Zustands kleiner als der erste ist, was z.B. bei Anwenden eines Operators geschehen kann, werden die Einträge getauscht und der Zustand erhält zusätzlich aufgrund der Antisymmetrie der fermionischen Wellenfunktion ein Vorzeichen.

In diesem Kapitel wird die Dynamik eines mikroskopischen Festkörpersystems, bestehend aus vier Gitterplätzen, im Hubbard-Modell untersucht.
Jeder Gitterplatz besitzt zwei Niveaus, die sich jeweils in der Spinmagnetquantenzahl $\sigma$ unterscheiden.
Auf jeweils vier Niveaus mit äquvalenter Spinausrichtung sind zwei Elektronen verteilt, sodass das System insgesamt halbgefüllt ist.
In Abbildung \ref{fig:system} ist eine Skizze des Systems mit beispielhafter Elektronenbesetzung abgebildet.

Die Elektronen können unter Beachtung des Pauliprinzips waagerecht oder senkrecht zwischen den Gitterplätzen unter Beibehaltung der Spinausrichtung springen.
Das System besteht daher effektiv aus zwei gekoppelten, spinlosen, halbgefüllten Festkörpersystemen mit jeweils vier Gitterplätzen, vier Niveaus.
In Kapitel \ref{sec:tbuntersuchung} wurde bereits herausgefunden, dass ein solches System sechs Zustände beinhaltet bzw. einen 6D-Hilbertraum umfasst.
Das gekoppelte System beinhaltet daher $6 \cdot 6=36$ Zustände bzw. umfasst einen 36D-Hilbertraum. Die Zustände des Gesamtsystems werden mit der Vorschrift
\begin{align}
  \ket{i} = \ket{f_i(\uparrow_1),f_i(\uparrow_2); f_i(\downarrow_1),f_i(\downarrow_2)}
  \label{eqn:hubbzustandsvorschrift}
\end{align}
dargestellt. Dabei ist $f_i(x_\alpha)$ eine Funktion, welche den Ort des Elektrons $x_\alpha$ im i-ten Zustand des Systems angibt.
In der Bezeichnung des Elektrons $x_\alpha$ stellt $x$ die Spinausrichtung und der Index $\alpha$ die Nummer des Elektrons dar, wobei die Elektronen eines Spinniveaus gegen den Uhrzeigersinn und beginnend beim Gitterplatz $1$ durchnummeriert werden.
Werden zwei Einträge einer Spinausrichtung in den Zuständen \eqref{eqn:hubbzustandsvorschrift} getauscht, so erhält der Zustand aufgrund der Antisymmetrie der fermionischen Wellenfunktion ein negatives Vorzeichen.
Die Einträge der $36$ Zustände aus Gleichung \eqref{eqn:hubbzustandsvorschrift} sind in Tabelle \ref{tab:hubbzustandsvorschrift} aufgelistet.

\begin{table}
  \centering
  \caption{.}
  \begin{minipage}[t]{0.35\linewidth}
    \begin{tabular}{S[table-format=2.0] S}
      \toprule
      {$i$} & {$\ket{i}$} \\
      \midrule
      1  & $\ket{1,2;1,2}$ \\
      2  & $\ket{1,2;1,3}$ \\
      3  & $\ket{1,2;1,4}$ \\
      4  & $\ket{1,2;2,3}$ \\
      5  & $\ket{1,2;2,4}$ \\
      6  & $\ket{1,2;3,4}$ \\
      7  & $\ket{1,3;1,2}$ \\
      8  & $\ket{1,3;1,3}$ \\
      9  & $\ket{1,3;1,4}$ \\
      10 & $\ket{1,3;2,3}$ \\
      11 & $\ket{1,3;2,4}$ \\
      12 & $\ket{1,3;3,4}$ \\
      13 & $\ket{1,4;1,2}$ \\
      14 & $\ket{1,4;1,3}$ \\
      15 & $\ket{1,4;1,4}$ \\
      16 & $\ket{1,4;2,3}$ \\
      17 & $\ket{1,4;2,4}$ \\
      18 & $\ket{1,4;3,4}$ \\
    \end{tabular}
  \end{minipage}
  \begin{minipage}[t]{0.35\linewidth}
    \begin{tabular}{S[table-format=2.0] S}
      \toprule
      {$i$} & {$\ket{i}$} \\
      \midrule
      19 & $\ket{2,3;1,2}$ \\
      20 & $\ket{2,3;1,3}$ \\
      21 & $\ket{2,3;1,4}$ \\
      22 & $\ket{2,3;2,3}$ \\
      23 & $\ket{2,3;2,4}$ \\
      24 & $\ket{2,3;3,4}$ \\
      25 & $\ket{2,4;1,2}$ \\
      26 & $\ket{2,4;1,3}$ \\
      27 & $\ket{2,4;1,4}$ \\
      28 & $\ket{2,4;2,3}$ \\
      29 & $\ket{2,4;2,4}$ \\
      30 & $\ket{2,4;3,4}$ \\
      31 & $\ket{3,4;1,2}$ \\
      32 & $\ket{3,4;1,3}$ \\
      33 & $\ket{3,4;1,4}$ \\
      34 & $\ket{3,4;2,3}$ \\
      35 & $\ket{3,4;2,4}$ \\
      36 & $\ket{3,4;3,4}$ \\
    \end{tabular}
  \end{minipage}
  \label{tab:hubbzustandsvorschrift}
\end{table}

Aus den Gleichungen \eqref{eqn:hamiltontb} und \eqref{eqn:hamiltonhubb} folgt der Hamiltonoperator für das Gesamtsystem
\begin{align}
  H_\text{hubb,sys} = -J \sum_{i=1}^4 \sum_{\sigma = \{ \downarrow,\uparrow \}} (c_{i+1,\sigma}^{\dag}c_{i,\sigma} + c_{i,\sigma}^{\dag}c_{i+1,\sigma}) + U\sum_{i=1}^4 \prod_{\sigma = \{ \downarrow,\uparrow \}} n_{i,\sigma},
  \label{eqn:hubbsys}
\end{align}
wobei die zusätzliche Bedingung periodischer Randbedingungen in den Summen über die Gitterplätze $i$ beachtet werden muss.
Der Hamilton-Operator wird durch die Vorschrift \eqref{eqn:matrixdarstellung} als $36 \times 36$ -Matrix dargestellt.
Dabei bestimmt der Vorfaktor $J$ des Tight-Binding-Terms \eqref{eqn:hubbsys} die Nichtdiagonalelemente und der
Vorfaktor $U$ des Elektronen-Potentialterms die Diagonalelemente der Hamilton-Matrix. Der Parameter $J$ wird im Folgenden auf $1$ gesetzt, sodass die Lösungen der Schrödingergleichung nur abhängig von $\frac{U}{J} = U$ sind.
Die berechnete Hamilton-Matrix ist in Tabelle \ref{tab:hubbhammatrix} abgebildet. Da der Hamiltonoperator $H$ nicht zeitabhängig ist, wird erneut die stationäre Schrödingergleichung \eqref{eqn:statschroedinger}
für das gegebene System aufgestellt und gelöst.
Zur Bestimmung der Eigenenergien und Eigenfunktionen von $H$ wird die Hamiltonmatrix für verschiedene $U/J$ mit Hilfe der Funktion // eig von octave // diagonalisiert. Die drei kleinsten Eigenwerte $\frac{E}{J}$ sind
in Tabelle \ref{tab:eigenwerte} für diverse $U/J$ aufsteigend aufgelistet und in Abbildung \ref{fig:eplot} gegen $U/J$ graphisch aufgetragen.

\begin{table}[h]
  \centering
  \caption{Datenpunkte der drei niedrigsten Eigenenergien des Festkörpersystems in Abhängigkeit vom Potential und in Einheiten von $J$.}
  \begin{tabular}{S[table-format=1.0] S[table-format=2.2] S[table-format=2.2] S[table-format=2.2]}
    \toprule
    {$U/J$} & {$E_0/J$} & {$E_1/J$} & {$E_2/J$} \\
    \midrule
    1  & -3.34 & -3.29 & -2.86 \\
    2  & -2.83 & -2.69 & -2.00 \\
    3  & -2.42 & -2.19 & -1.42 \\
    4  & -2.10 & -1.81 & -1.07 \\
    5  & -1.84 & -1.51 & -0.84 \\
    6  & -1.63 & -1.29 & -0.70 \\
    7  & -1.46 & -1.12 & -0.59 \\
    8  & -1.32 & -0.99 & -0.51 \\
    9  & -1.20 & -0.88 & -0.45 \\
    10 & -1.10 & -0.80 & -0.41 \\
    \bottomrule
  \end{tabular}
  \label{tab:eigenwerte}
\end{table}

\begin{figure}
  \centering
  \includegraphics{build/Hubb_eplot.pdf}
  \caption{plot.}
  \label{fig:eplot}
\end{figure}

Für ein verschwindendes Potential $U$ nähern sich die drei niedrigsten Eigenenergien dem Wert $E/J = -4$ an.
Das ist dadurch erklärbar, dass sich das 8-Niveau-Festkörpersystem im Grenzfall $U/J = 0$ wie zwei gekoppelte, nicht miteinander wechselwirkende 4-Niveau-Festkörpersysteme verhält.
Die niedrigste Eigenenergie $E/J = -2$ des 4-Niveau-Festkörpersystems   ist zweifach entartet.
Daher gibt es vier Möglichkeiten


\appendix
% Hier beginnt der Anhang, nummeriert in lateinischen Buchstaben
\chapter{Ein Anhangskapitel}

Hier könnte ein Anhang stehen, falls Sie z.B. Code, Konstruktionszeichnungen oder Ähnliches mit in die Arbeit bringen wollen. Im Normalfall stehen jedoch alle Ihre Resultate im Hauptteil der Bachelorarbeit und ein Anhang ist überflüssig.


\backmatter
\printbibliography

\cleardoublepage
\thispagestyle{empty}
\section*{Eidesstattliche Versicherung}
Ich versichere hiermit an Eides statt, dass ich die vorliegende Abschlussarbeit mit dem Titel \enquote{\thetitle} selbstständig und ohne unzulässige fremde Hilfe erbracht habe.
Ich habe keine anderen als die angegebenen Quellen und Hilfsmittel benutzt, sowie wörtliche und sinngemäße Zitate kenntlich gemacht. 
Die Arbeit hat in gleicher oder ähnlicher Form noch keiner Prüfungsbehörde vorgelegen.

\vspace*{1cm}\noindent
\begin{center}
  \begin{tabular}{@{}p{0.4\textwidth}@{\hspace{0.15\textwidth}}p{0.4\textwidth}@{}}
  \rule{\linewidth}{0.25pt}& \rule{\linewidth}{0.25pt}\\
  Ort, Datum & Unterschrift
  \end{tabular}
\end{center}

\subsection*{Belehrung}
Wer vorsätzlich gegen eine die Täuschung über Prüfungsleistungen betreffende Regelung einer Hochschulprüfungsordnung verstößt, handelt ordnungswidrig.
Die Ordnungswidrigkeit kann mit einer Geldbuße von bis zu \SI[round-mode=places, round-precision=2]{50000}{€} geahndet werden. 
Zuständige Verwaltungsbehörde für die Verfolgung und Ahndung von Ordnungswidrigkeiten ist der Kanzler/die Kanzlerin der Technischen Universität Dortmund. 
Im Falle eines mehrfachen oder sonstigen schwerwiegenden Täuschungsversuches kann der Prüfling zudem exmatrikuliert werden \mbox{(\S\,63 Abs. 5 Hochschulgesetz --HG--).}

Die Abgabe einer falschen Versicherung an Eides statt wird mit Freiheitsstrafe bis zu 3 Jahren oder mit Geldstrafe bestraft.

Die Technische Universität Dortmund wird ggf.\ elektronische Vergleichswerkzeuge (wie z.\,B.\ die Software \enquote{turnitin}) zur Überprüfung von Ordnungswidrigkeiten in Prüfungsverfahren nutzen. \\[\baselineskip]

\noindent Die oben stehende Belehrung habe ich zur Kenntnis genommen.\\[1cm]
\begin{center}
\begin{tabular}{@{}p{0.4\textwidth}@{\hspace{0.15\textwidth}}p{0.4\textwidth}@{}}
\rule{\linewidth}{0.25pt}& \rule{\linewidth}{0.25pt}\\
Ort, Datum & Unterschrift
\end{tabular}
\end{center}

\end{document}
